\documentclass[10pt,a4paper]{amsart}
\usepackage[margin=19mm]{geometry}
\usepackage{amsmath,amssymb,mathtools}
\usepackage[scr=rsfs,cal=euler]{mathalfa}
\usepackage{xfrac}
\usepackage[unicode,hidelinks,bookmarks=false]{hyperref}
\newcommand{\ie}{i.e.\ }
\newcommand{\cf}{cf.\ }
\newcommand{\resp}{respectively\ }
\newcommand{\Subsection}{\subsection}
\providecommand{\nobreakdash}{\nobreak\hbox{-}}
\setlength{\emergencystretch}{2em}
\multlinegap0pt
\usepackage{bm}
\usepackage{bbm}
\usepackage{dsfont}
\usepackage{xspace}
\usepackage{cancel}
\usepackage{mathtools}
\usepackage{amsthm}
\makeatletter
\@ifundefined{theorem}{\newtheorem{theorem}{Theorem}}{}
\@ifundefined{definition}{\newtheorem{definition}[theorem]{Definition}}{}
\@ifundefined{proposition}{\newtheorem{proposition}[theorem]{Proposition}}{}
\makeatother
\newcommand{\hC}{\widehat{\mathbb{C}}}
\newcommand{\mD}{\mathbb{D}}
\newcommand{\mZ}{\mathbb{Z}}
\renewcommand {\tilde} {\widetilde}
\title[Rational lemniscates and the matching problem]{Rational lemniscates and the matching problem}
\author{Kirill Lazebnik, Pierre-Olivier Paris\'e, Malik Younsi}
\date{Source: arXiv:2507.03547v1 (CC BY 4.0)}
\begin{document}
\maketitle

\noindent\emph{Source and licence.} Rational lemniscates and the matching problem. Version arXiv:2507.03547v1 (CC BY 4.0), distributed under the Creative Commons Attribution 4.0 International licence, as recorded in the arXiv licence metadata for this exact version and shown on the abstract page https://arxiv.org/abs/2507.03547v1. Full rights notice: licenses/arxiv-2507.03547-CC-BY-4.0.txt.\par\smallskip

\noindent\textit{Editorial scope.} Counted unit: the Proposition whose source label is \texttt{proper2} in the article (source0.tex lines 451--460, source label \texttt{proper2}), delivered together with the article's own complete original proof of it (source0.tex lines 462--499, one \texttt{proof} environment of 38 lines). No mathematics is written, repaired or re-derived: statement and proof are the article's own text line for line. Required context retained uncounted: lemniscatedomain (lines 379--382); extension (lines 391--398); proper1 (lines 433--439). Every definition, display and label of those lines is retained unchanged. Omitted from this package and not redistributed: 1--378, 383--390, 399--432, 440--450, 461--461, 500--966. Nothing inside a retained excerpt is omitted or altered: every retained body is delivered exactly as the article's lines are written, so no omission receipt of any kind accompanies this package. Citations: the retained excerpts carry no citation command at all, so no bibliography entry of the article is transcribed and no reference is redistributed. Reference adaptation: none was needed and none was made; every reference inside the delivered unit resolves inside it, so no unresolved reference or citation is produced. Printed slips kept literal: no printed word, symbol, hypothesis or number of the counted statement or of its proof was altered. Layout and class: the article's own class is \texttt{amsart} with packages including hyperref, amssymb, amsmath, amsthm, enumerate, graphicx, wrapfig, color, subcaption, float, tikz; the wrapper is the lane's pinned \texttt{amsart} editorial wrapper at 10pt on A4, which restates the theorem-like environments the retained text needs and adds the article's own macro definitions that the retained text needs (\texttt{\textbackslash hC}, \texttt{\textbackslash mD}, \texttt{\textbackslash mZ}, \texttt{\textbackslash tilde}), with extra wrapper packages (bm, bbm, dsfont, xspace, cancel, mathtools). The delivered numbering is the unit's own. Assets: no retained span contains a figure environment, a graphics-inclusion command, a PDF-page inclusion, a table or a TikZ picture; no image file is copied into this package, the package declares no asset dependency and no image file is redistributed with it. Licence: version arXiv:2507.03547v1 of this article is distributed under the Creative Commons Attribution 4.0 International licence, as recorded in the arXiv licence metadata for this exact version and shown on the abstract page https://arxiv.org/abs/2507.03547v1. A full rights notice is delivered at licenses/arxiv-2507.03547-CC-BY-4.0.txt. Characters: every retained excerpt is pure ASCII, so the delivered engine, XeLaTeX with its default Latin Modern text font, reports no missing character.
\begin{definition}
\label{lemniscatedomain}
A domain $\Omega \subset \hC$ is called a \textit{lemniscate domain} if its boundary $\partial \Omega$ is a lemniscate graph. In the special case where $\partial \Omega$ consists of finitely many pairwise disjoint analytic Jordan curves, then we call $\Omega$ an \textit{analytic Jordan domain}.
\end{definition}

\begin{proposition}
\label{extension}
Let $\Omega$ be a lemniscate domain. Suppose that every edge of $\partial \Omega$ is analytic. Then the Green's function for $\Omega$ extends to be harmonic on a neighborhood of $\partial \Omega \setminus V$, where $V \subset \partial \Omega$ is the set of vertices. Moreover we have, for $w \in \Omega$,
$$-\frac{\partial g_\Omega (\zeta , w )}{\partial n_\zeta} > 0 \qquad (\zeta \in \partial \Omega \setminus V),$$
where $n_\zeta$ is the unit outer normal at $\zeta$, and
$$d\omega(\zeta,w,\Omega) = -\frac{1}{2\pi} \frac{\partial g_\Omega (\zeta , w)}{\partial n_\zeta} |d\zeta|$$
as measures on $\partial \Omega \setminus V$.
\end{proposition}

\begin{proposition}
\label{proper1}
Let $\Omega \subset \hC$ be a lemniscate domain and let $B:\Omega \to \mD^*$ be a proper holomorphic mapping. If $w_1, \dots, w_n \in \Omega$ are the poles of $B$ repeated according to multiplicities, then
$$
\log{|B(z)|} = \sum_{k=1}^n g_{\Omega}(z,w_k) \qquad (z \in \Omega).
$$
\end{proposition}

\begin{proposition}
\label{proper2}
Let $\Omega \subset \hC$ be a lemniscate domain with boundary components $\Gamma_1, \dots, \Gamma_N$, and let $w_1, \dots, w_n$ be points in $\Omega$, not necessarily distinct. Then there exists a proper holomorphic mapping $B:\Omega \to \mD^*$ with poles precisely at $w_1, \dots, w_n$ if and only if for all $j \in \{1, \dots, N\}$ we have
\begin{equation}
\label{eqhm}
\sum_{k=1}^n \omega(\Gamma_j, w_k, \Omega) \in \mZ.
\end{equation}


\end{proposition}

\begin{proof}
We may assume $\Omega$ is bounded.

Suppose that Equation (\ref{eqhm}) holds, and first assume that $\Omega$ is an analytic Jordan domain. Consider the function
$$
v(z):=\sum_{k=1}^n g_{\Omega}(z,w_k) \qquad (z \in \Omega).
$$
Then $v$ is harmonic on $\Omega \setminus \{w_1, \dots, w_n\}$ and has a harmonic extension to a neighborhood of $\partial \Omega$, by Proposition \ref{extension}. The harmonic conjugate $\tilde{v}(z)$ is well-defined locally, but when we analytically continue along a given boundary curve $\Gamma_j$, the value of $\tilde{v}(z)$ changes by
$$
\int_{\Gamma_j} \frac{\partial v}{\partial n_\zeta} \, |d\zeta|.
$$
Using Equation (\ref{eqhm}) and Proposition \ref{extension}, it easily follows that the value of $\tilde{v}(z)$ changes by an integer multiple of $2\pi$ under harmonic continuation. Thus
$$
B(z):=\exp{\left(v(z) + i \tilde{v}(z)\right)} \qquad (z \in \Omega)
$$
is a well-defined holomorphic function, and it is readily checked that $B$ is a proper holomorphic mapping of $\Omega$ onto $\mD^*$ with poles precisely at $w_1, \dots, w_n$. This proves the desired implication, but only in the special case where $\Omega$ is an analytic Jordan domain. Now, if $\Omega$ is a general lemniscate domain, then consider a conformal map $f:\Omega \to \Omega'$ onto an analytic Jordan domain $\Omega'$, as in the remark following Definition \ref{lemniscatedomain}. Then $f$ has a continuous and injective extension to $\overline{\Omega} \setminus V$, where $V \subset \partial \Omega$ is the set of vertices. Since $V$ is a finite set and therefore has zero harmonic measure, it follows from Equation (\ref{eqhm}) and the conformal invariance of harmonic measure that
$$
\sum_{k=1}^n \omega(\Gamma', f(w_k), \Omega') \in \mZ
$$
for every boundary component $\Gamma'$ of $\Omega'$. By the analytic case, there exists a proper holomorphic mapping $\tilde{B}:\Omega' \to \mD^*$ with poles precisely at $f(w_1), \dots, f(w_n)$. Then $B:=\tilde{B} \circ f$ gives a proper holomorphic mapping of $\Omega$ onto $\mD^*$ with poles precisely at $w_1, \dots, w_n$, as required.
\\

Conversely, suppose that there exists a mapping $B$ with the specified properties. Composing with a conformal map if necessary, we may assume once again that $\Omega$ is an analytic Jordan domain. Recall that if $\Gamma$ is a boundary component of $\Omega$ and if $h$ is a smooth function defined in a neighborhood of $\Gamma$, then the \textit{period} of $h$ along $\Gamma$ is
$$
\mathrm{per} (h , \Gamma ) := \int_{\Gamma} \frac{\partial h}{\partial n_\zeta}\, |d\zeta| .
$$
Now, by Proposition \ref{proper1}, we have
$$
\log{|B(z)|} = \sum_{k=1}^n g_{\Omega}(z,w_k) \qquad (z \in \Omega).
$$
But $B$ is single-valued, so that for each $j \in \{1, \dots, N\}$, we have
\begin{align*}
\sum_{k=1}^n \omega(\Gamma_j,w_k,\Omega) &= \sum_{k=1}^n \left(- \frac{1}{2\pi} \operatorname{per}(g_{\Omega}(z,w_k), \Gamma_j)\right)\\
&= - \frac{1}{2\pi} \operatorname{per}\left(\sum_{k=1}^n g_{\Omega}(z,w_k), \Gamma_j\right)\\
&= - \frac{1}{2\pi} \operatorname{per}(\log{|B(z)|}, \Gamma_j) \in \mathbb{Z},
\end{align*}
as required.
\end{proof}

\end{document}
